\section{Dokumentation der Entwicklungspipeline}
\label{sec:ergebnisse}

Damit die entwickelten Digitalen Zwillinge über die Laufzeit dieser Arbeit hinaus im LFP genutzt, gewartet und auf weitere Anlagen übertragen werden können, mussten Entwicklungsergebnisse, Reproduktionsschritte und Übergabeartefakte durchgängig dokumentiert werden. Dieses Kapitel fasst zunächst die beiden entstandenen Softwareerweiterungen zusammen, beschreibt anschließend die Reproduzierbarkeitsdokumentation und schließt mit einer Gegenüberstellung beider Anlagen.

\subsection{Omniverse-Extension und ROS\,2-Anbindung als Entwicklungsergebnisse}

Im Rahmen dieser Arbeit wurden zwei funktionsfähige Softwareerweiterungen für NVIDIA Isaac Sim entwickelt, erprobt und an das LFP übergeben.

\textbf{Maze-Runner-Extension:} Die entwickelte modulare Extension im \emph{Omniverse Kit Extension Framework} gliedert sich in zwölf spezialisierte Module (vgl.\ Abschnitt~\ref{subsubsec:mazerunner_extension}) und koppelt die Siemens-SPS S7-1200 über den \emph{DigitalTwinService} bidirektional an den Digitalen Zwilling in NVIDIA Isaac Sim. Die datengetriebene Konfiguration über \texttt{nodes\_db.json} ermöglicht die Übertragung der Extension auf weitere SPS-gestützte Anlagen ohne Quellcodeänderungen. Der Steuerpfad von der Extension zur SPS erreicht dabei eine Latenz zwischen \SI{100}{\milli\second} und \SI{500}{\milli\second} (vgl.\ Abschnitt~\ref{subsubsec:mazerunner_sps}). Die Rückmeldung der Sensorwerte aus der SPS an den Digitalen Zwilling über die MQTT-WebSocket-Schnittstelle ließ sich im Laborbetrieb dagegen nicht durchgängig zuverlässig nachweisen, worauf Abschnitt~\ref{subsec:gesamtbewertung} genauer eingeht.

\textbf{Dobot-Magician-Extension:} Für den Dobot Magician entstand mit \texttt{Dobot\_DT\_IsaacSim\_Extension} eine eigenständige Kit-Extension, deren Klasse \texttt{DobotROSNode} über die \emph{pydobot}-Bibliothek zyklisch die Pose des realen Roboters ausliest und die berechneten Gelenkwinkel direkt über die in Isaac Sim integrierte Dynamic-Control-API auf die Articulation des importierten Robotermodells schreibt. In Gegenrichtung sendet die Extension über ein Bedienpanel Jog-Befehle, Vakuumgreifer-Schaltungen sowie automatisierte Pick-and-Place-Sequenzen an den realen Gelenkarmroboter, wobei Letztere zuvor per Teach-In am physischen Gelenkarmroboter aufgezeichnet wurden. Beide Erweiterungen sind im GitHub-Repository versioniert und gemäß der Reproduzierbarkeitsdokumentation (vgl.\ Abschnitt~\ref{subsec:reprodoku}) nachvollziehbar.

\subsection{Reproduzierbarkeitsdokumentation}
\label{subsec:reprodoku}

Damit beide Digitalen Zwillinge unabhängig von den beteiligten Personen reproduziert werden können, wurde für jede Anlage eine eigenständige Anleitung erstellt und in komplementären Medienformaten an das LFP übergeben.

\subsubsection{Methodische Schritt-für-Schritt-Anleitung -- Maze Runner}
Aus den Erfahrungen der Modellerstellung wurde eine generische Anleitung in sieben Schritten abgeleitet (vgl.\ Anhang~\ref{anhang:repro}):
\begin{enumerate}
    \item Anforderungsanalyse und Schnittstellenmatrix der Zielanlage,
    \item CAD-Aufbereitung und Export nach OpenUSD,
    \item Definition kinematischer Strukturen in NVIDIA Isaac Sim (Articulations),
    \item Auswahl und Einrichtung der geeigneten Middleware (OPC~UA, MQTT, ROS\,2),
    \item Implementierung der bidirektionalen Datenkopplung,
    \item Validierung anhand definierter Testszenarien,
    \item Übergabe und Versionierung im GitHub-Repository.
\end{enumerate}

\subsubsection{Reproduzierbarkeitsdokumentation Dobot Magician}
Analog zur Maze-Runner-Dokumentation wurde für den Dobot Magician eine eigenständige Inbetriebnahme-Anleitung erstellt. Diese umfasst:
\begin{enumerate}
    \item Einrichtung von \emph{pydobot} unter Windows sowie Installation der Extension \texttt{Dobot\_DT\_IsaacSim\_Extension} über den Extension-Manager oder das Hilfsskript \texttt{install\_extension.py},
    \item URDF-Import in NVIDIA Isaac Sim mit den erforderlichen Parametereinstellungen (\emph{Merge Fixed Joints}, \emph{Fix Base Link}, deaktiviertes \emph{Instanceable}),
    \item Verifikation der Dynamic-Control-Kopplung anhand der Gelenkzustandsübertragung,
    \item Betrieb des Software-Mocks für hardware-unabhängige Entwicklung und Tests.
\end{enumerate}
Alle Schritte sind im GitHub-Repository versioniert. Die Anleitung ist so strukturiert, dass sie ohne physischen Roboter vollständig nachvollzogen werden kann.

\subsubsection{Medienformate}
Die Anleitungen wurden in komplementären Formaten aufbereitet:
\begin{enumerate}
    \item \textbf{Foliensatz} (PowerPoint, 32 Folien) als didaktische Grundlage für Lehrveranstaltungen,
    \item \textbf{VR-Demonstration} auf Meta Quest~3 zur immersiven Inspektion der Digitalen Zwillinge,
    \item \textbf{Maze-Runner-GitHub-Repository} mit allen im Rahmen dieser Arbeit erarbeiteten Dateien, von der Installation von NVIDIA Isaac Sim über die Konfiguration der SPS und die Anbindung an die Netzwerkinfrastruktur der Hochschule RheinMain bis hin zur Extension selbst mit ihren Modulen sowie der Maze-Runner-USD-Szene,
    \item \textbf{Dobot-Extension-GitHub-Repository} mit der technischen Dokumentation der Dobot-Magician-Extension, der URDF-Datei mit den zugehörigen CAD-Dateien sowie Hilfsskripten, etwa zur Kalibrierung, zur Installation der Extension oder zum Setzen des Extension-Pfads unter Windows.
\end{enumerate}
Beide Anleitungen wurden von mehreren Personen getestet, deren Rückmeldungen eingearbeitet wurden.

\subsubsection{Übergabe an das LFP}
Alle erstellten Artefakte -- USD-Stages, URDF-Dateien, Python-Extensions, OPC-UA-Konfigurationen, Reproduktionsanleitungen und die schriftliche Ausarbeitung -- wurden gemäß Aufgabenstellung digital an das LFP übergeben und in einem GitHub-Repository versioniert. Eine Veröffentlichung innerhalb der für IEEE geplanten Schriftenreihe des LFP ist vorgesehen.

\subsection{Gegenüberstellung und Gesamtbewertung}
\label{subsec:gesamtbewertung}

Beide Anlagen bestätigen die Machbarkeit einer durchgängigen Datenkette von CAD über OpenUSD bis zur Echtzeit-Simulation in NVIDIA Isaac Sim ohne proprietäre Konvertierungsschritte~\cite{pixar_openusd}. Die Protokollwahl, OPC~UA und MQTT für die SPS-Anlage sowie ROS\,2 für den Roboter, richtet sich nach den unterschiedlichen Kommunikationsanforderungen beider Anlagen.

Beim Maze Runner ließ sich die Rückmeldung der Sensorwerte aus der SPS an den Digitalen Zwilling über die MQTT-WebSocket-Schnittstelle im Laborbetrieb nicht durchgängig echtzeitfähig nachweisen. Über den eingerichteten Subscribe-Mechanismus gingen zwar MQTT-Pakete ein, jedoch nicht durchgängig echtzeitfähig und zuverlässig genug, um als Grundlage für einen echtzeitfähigen Digitalen Zwilling zu dienen. Als mögliche Ursache kommt weniger eine interne Netzwerktrennung im Backend als vielmehr die Netzwerktrennung zwischen dem Maze-Runner-Netzwerk und dem Backend der Hochschule RheinMain über die bestehende VPN-Verbindung infrage, ohne dass die dortigen IT-Verantwortlichen hierfür eine abschließende Erklärung liefern konnten. Die Kopplung erreicht damit im praktischen Betrieb nicht durchgängig die von Kritzinger et al.\ definierte Stufe eines bidirektionalen Digitalen Zwillings~\cite{kritzinger2018}. Da der Steuerpfad von der Simulation zur SPS zuverlässig funktioniert, während die Rückmeldung eingeschränkt bleibt, entspricht der erreichte Zustand eher einer erweiterten Fernsteuerung mit ergänzendem, nicht durchgängig verfügbarem Digitalen Schatten.

Beim Dobot Magician ist zumindest die Richtung vom realen Gelenkarmroboter zur Simulation durchgängig automatisiert. Gelenkwinkel werden über die Dynamic-Control-API kontinuierlich ausgelesen und auf die Simulation übertragen, während Steuerbefehle und aufgezeichnete Pick-and-Place-Sequenzen in der Gegenrichtung über die UI-Buttons des Bedienpanels der Extension manuell ausgelöst werden. Für die Aufzeichnung dieser Sequenzen wird am realen Gelenkarmroboter die Motorbremse gelöst, sodass der Arm per Hand in die gewünschte Position geführt und die Pose per Teach-In übernommen werden kann. Hinsichtlich der Skalierbarkeit wächst die GPU-Last linear mit der Anzahl gleichzeitig simulierter Articulations, was die in Kapitel~\ref{sec:grundlagen} adressierte Infrastrukturanforderung bestätigt. Für vergleichbare Laboraufbauten wird eine Kombination beider Ansätze empfohlen, OPC~UA und MQTT für die Anlagenebene sowie ROS\,2 für die Roboterebene, wobei die Zuverlässigkeit der Rückmeldekanäle vor einer produktiven Nutzung gesondert zu prüfen ist.

% =====================================================================
% GEPARKT (2026-07-05): Abschnitt "Projektorganisation" aus Kapitel 4.1
% entfernt, damit der Text nicht verloren geht. Bei Bedarf hier wieder
% einordnen und als \subsection aktivieren.
% =====================================================================
\iffalse
\subsection{Projektorganisation}
Zunächst entstand im Rahmen des IIM-Projektes der Digitale Zwilling des Maze-Runner,
wobei die erlangten Kenntnisse festgehalten wurden und als Grundlage für die weitere Arbeit dienten.
Der Maze-Runner lag zu Projektbeginn weder vollständig aufgebaut noch betriebsbereit vor.
Von der HSRM standen eine CAD-Baugruppe der Anlage, Dokumentation
vergangener Projekte mit Stücklisten sowie auf Moodle hinterlegte Trainingslevel zur
Einarbeitung in den Anlagenbau- und SPS-Programmierung zur Verfügung.
Die CAD-Baugruppe bildete dabei die Grundlage für das digitale Modell.
Während für die Inbetriebnahme des Maze-Runners die Bereiche Mechanik, Konstruktion, Pneumatik, Elektrik, Steuerungstechnik und Netzwerktechnik relevant waren, betrafen den Digitalen Zwilling vorrangig die Schnittstellen zur Steuerungstechnik, Softwareentwicklung sowie die Einarbeitung in die Simulationsumgebung NVIDIA Isaac Sim.
Softwareentwicklung, Simulation und Dokumentation wurden im Homeoffice bearbeitet.
Parallel dazu wurden die ortsgebundenen Aufgaben an der HsH durchgeführt, wobei
kurze Rücksprachen mit den betreuenden Personen im Labor zur Festhaltung des
aktuellen Fortschritts sowie zur Diskussion weiterer Anforderungen genutzt wurden.
Bei steuerungstechnischen Fragen konnte auf Kollegen des Fachbereichs der HsH zurückgegriffen werden.
Wartezeiten, die etwa durch laufende Inbetriebnahmeschritte oder das Eintreffen
bestellter Teile entstanden, wurden genutzt, um anderen Bereichen nachzugehen und
dem Ziel eines Digitalen Zwillings schrittweise näherzukommen.

Anlagenteile, die zum Zeitpunkt der Modellierung noch nicht in Betrieb waren, wurden
im Digitalen Zwilling anhand der Dokumentation der HSRM virtuell in
Betrieb genommen und später an der realen Anlage verifiziert oder angepasst.

Nach erfolgreichem Abschluss des IIM-Projektes wurde im Rahmen der Bachelorarbeit
ein Digitaler Zwilling des Dobot Magician erstellt.
Da kein CAD-Modell vorhanden war, musste eines direkt beim Hersteller angefragt werden.
Für den Dobot Magician wurde ebenfalls eine Extension entwickelt, wobei die Erstellungszeit
durch die Kenntnisse aus der Einarbeitung in NVIDIA Isaac Sim und der vorherigen Extension-Entwicklung
erheblich verkürzt werden konnte.
Die Schritt-für-Schritt-Reproduzierbarkeitsdokumentation des Maze-Runner wird durch die technische Dokumentation
des Digitalen Zwillings des Dobot Magician so ergänzt, dass sich beide Anlagen
reproduzieren lassen.
\fi
